\documentclass[10pt,a4paper]{amsart}
\usepackage[margin=19mm]{geometry}
\usepackage{amsmath,amssymb,mathtools}
\usepackage[scr=rsfs,cal=euler]{mathalfa}
\usepackage{xfrac}
\usepackage[unicode,hidelinks,bookmarks=false]{hyperref}
\newcommand{\ie}{i.e.\ }
\newcommand{\cf}{cf.\ }
\newcommand{\resp}{respectively\ }
\newcommand{\Subsection}{\subsection}
\providecommand{\nobreakdash}{\nobreak\hbox{-}}
\setlength{\emergencystretch}{2em}
\multlinegap0pt
\usepackage{amssymb}
\usepackage{mathtools}
\usepackage[shortlabels]{enumitem}
\newtheorem{lemma}{Lemma}
\newcommand{\R}{\mathbb R}
\title{Classification of compact Lagrangian self-similar submanifolds with Legendrian capillary boundary in the unit ball}
\author{Dong Gao, Yong Luo, Hui Ma, Jiabin Yin}
\date{Source: arXiv:2608.14431v1 (CC BY 4.0)}
\begin{document}
\maketitle

\noindent\emph{Source and licence.} Classification of compact Lagrangian self-similar submanifolds with Legendrian capillary boundary in the unit ball. Version arXiv:2608.14431v1 (CC BY 4.0), distributed by arXiv under the Creative Commons Attribution 4.0 International licence, http://creativecommons.org/licenses/by/4.0/, as recorded in the arXiv licence metadata for this exact version and shown on the abstract page https://arxiv.org/abs/2608.14431v1. Full licence notice: \texttt{licenses/arxiv-2608.14431-CC-BY-4.0.txt}.\par\smallskip

\noindent\textit{Editorial scope.} Counted unit: the article's lemma lem:cauchy (source0.tex lines 691--729). The article's complete original proof of it is delivered with it (source0.tex lines 731--853, one \texttt{proof} environment of 123 lines). No mathematics is written, repaired or re-derived: the statement and the proof are the article's own lines, in the article's own notation, and no retained line is substituted or re-flowed.  Retained uncounted as the source-local context the proof needs: lines 663--670; lines 665--668.  Nothing was omitted from any retained span, so the package carries no omission record.  No retained line contains a citation, so the wrapper carries no bibliography and no bibliography entry.  No retained line contains a figure environment, an image-inclusion command or drawing code, so no image is delivered and the package declares no asset dependency.  Editorial formatting only, in addition: an AMS \texttt{amsart} preamble that restates the article's own declarations and macros used by the retained lines, namely the environments \texttt{lemma} on separate counters, together with the standard packages those lines need.  The retained environments print in this document as Lemma 1 in order of appearance, whereas the article prints the same items with the numbering of its own full document.  The complete native source of the article as distributed in the arXiv e-print for this exact version is bundled unchanged as \texttt{sources/207669/source0.tex}.
\begin{lemma}\label{lem:ucp-external}
Let $w\in W^{2,2}_{\mathrm{loc}}(\Omega;\R^m)$, where $\Omega\subset\R^n$ is connected.  Suppose
\begin{equation}\label{eq:ucp-inequality}
        |A^{ij}\partial_{ij}w|
        \le C\bigl(|w|+|Dw|\bigr)
\end{equation}
almost everywhere, where $A^{ij}=A^{ji}$ are Lipschitz and uniformly elliptic and the same scalar operator $A^{ij}\partial_{ij}$ acts on each component.  If $w$ vanishes on a nonempty open subset of $\Omega$, then $w\equiv0$ in $\Omega$.
\end{lemma}

\begin{equation}
        |A^{ij}\partial_{ij}w|
        \le C\bigl(|w|+|Dw|\bigr)
\end{equation}

\begin{lemma}\label{lem:cauchy}
Fix $\varepsilon\in\{-1,0,1\}$.  For $a=1,2$, let
$$
        X_a:M_a^n\longrightarrow\mathbb R^N
$$
be smooth immersed solutions of
$$
        H_a+\varepsilon X_a^\perp=0.
$$
Let $\Sigma_a$ be compact boundary components of $M_a$, let
$$
        F:\Sigma_1\longrightarrow\Sigma_2
$$
be a diffeomorphism, and let $\eta_a$ denote the inward unit conormal
along $\Sigma_a$.  Suppose that, for every $p\in\Sigma_1$,
$$
        X_1(p)=X_2(F(p))
$$
and
$$
        dX_1|_p\bigl(\eta_1(p)\bigr)
        =
        dX_2|_{F(p)}\bigl(\eta_2(F(p))\bigr).
$$
Then there exist one-sided neighborhoods $V_a$ of $\Sigma_a$ and a
diffeomorphism
$$
        \Phi:V_1\longrightarrow V_2
$$
such that
$$
        \Phi|_{\Sigma_1}=F,
        \qquad
        X_1=X_2\circ\Phi
        \quad\text{on }V_1.
$$


\end{lemma}

\begin{proof}
We work sheetwise near a fixed point $p_1\in\Sigma_1$ and its corresponding
point $p_2=F(p_1)\in\Sigma_2$.  The possible presence of other preimages of
the same ambient point is irrelevant, because an immersion is an embedding on
a sufficiently small neighborhood of each prescribed preimage.  Shrink the
two neighborhoods so that both chosen sheets are embedded.

Let $P$ be their common tangent $n$-plane at the boundary point, and let
$\pi_P$ be orthogonal projection onto $P$.  The differential of
$\pi_P\circ X_a$ is an isomorphism at $p_a$; hence, after shrinking again,
each sheet is a graph over a one-sided domain in $P$.  The equality of the
boundary parametrizations implies that the projected boundary hypersurfaces
are the same.  Equality of the specified inward conormals implies that the two
projected sheets lie on the same side of that hypersurface.  We may therefore
choose one connected one-sided neighborhood $\Omega^+\subset P$, contained in
both projected domains, and write
\begin{equation}
        Y_a(y)=y+u_a(y),
        \qquad
        u_a:\Omega^+\longrightarrow P^\perp,
        \qquad a=1,2.
\end{equation}
Here the origin and a fixed translation have been absorbed into the affine
coordinate $y$.

Let $\Gamma=\partial\Omega^+$ denote the common projected boundary portion.
On $\Gamma$, equality of the boundary immersions gives $u_1=u_2$.  For a graph
over a fixed plane, its tangent space is the graph of the unique linear map
$Du_a:P\to P^\perp$.  Since the full tangent $n$-planes of the two immersed
sheets agree at every corresponding boundary point, it follows that
$Du_1=Du_2$ on $\Gamma$.  The conormal hypothesis is used here to select the
same one-sided graph; it rules out comparing the upper sheet of one immersion
with the lower sheet of the other.  Thus, with $w=u_1-u_2$,
\begin{equation}
        w=0,
        \qquad
        Dw=0
        \qquad\ \text{on }\Gamma.
\end{equation}

Choose a smooth diffeomorphism
$\chi:B_r^+\to\Omega^+$ which maps $I_r=\{y_n=0\}$ onto $\Gamma$, and replace
$u_a$ by $u_a\circ\chi$.  Because the same fixed change of independent
variables is used for both systems, the transformed principal operator remains
a scalar uniformly elliptic operator acting identically on every normal
component.  More explicitly, the affine graph equation
\begin{equation*}
        g^{ij}(Du_a)\partial_{ij}u_a^\alpha
        +\varepsilon\langle y+u_a,N_\alpha(Du_a)\rangle=0
\end{equation*}
transforms into a quasilinear system whose leading coefficient is
\begin{equation*}
        \widehat g_a^{kl}
        =g_a^{ij}\,\partial_i(\chi^{-1})^k\,\partial_j(\chi^{-1})^l,
\end{equation*}
while the second derivatives of $\chi$ occur only in lower-order terms.
After this common flattening, we have
\begin{equation}\label{eq:full-cauchy-data}
        w=0,
        \qquad
        Dw=0
        \qquad\ \text{on }I_r.
\end{equation}

Subtract the two transformed quasilinear systems.  For the leading terms,
use the fundamental theorem of calculus along the segment
$u_2+t(u_1-u_2)$ in the variables $(u,Du)$.  The result is a linear system
\begin{equation}
        A^{ij}(y)\partial_{ij}w^\alpha
        +B^i_{\alpha\beta}(y)\partial_iw^\beta
        +C_{\alpha\beta}(y)w^\beta=0,
\end{equation}
where $A^{ij}=A^{ji}$ is smooth, uniformly elliptic, and independent of the
normal index $\alpha$.  Consequently,
\begin{equation}
        |A^{ij}\partial_{ij}w|
        \le C\bigl(|w|+|Dw|\bigr)
        \qquad\text{in }B_r^+.
\end{equation}

Define the zero extension
\begin{equation}
        \widetilde w(y)=
        \begin{cases}
        w(y),&y_n\ge0,\\
        0,&y_n<0.
        \end{cases}
\end{equation}
We verify the required Sobolev regularity.  For a scalar component and a test
function $\zeta$, integration by parts in the upper half-ball gives
\begin{equation}
        \langle\partial_i\widetilde w,\zeta\rangle
        =\int_{B_r^+}(\partial_iw)\zeta,
\end{equation}
because the boundary term contains the trace of $w$, which is zero.  A second
integration by parts gives
\begin{equation}\label{eq:second-extension-derivative}
        \langle\partial_{ij}\widetilde w,\zeta\rangle
        =\int_{B_r^+}(\partial_{ij}w)\zeta,
\end{equation}
because the remaining interface term contains the trace of $\partial_jw$,
which also vanishes.  Hence $\widetilde w\in W^{2,2}_{\mathrm{loc}}(B_r)$ and
its weak second derivatives are precisely the zero extensions of those of
$w$.

Extend $A^{ij}$ to Lipschitz uniformly elliptic coefficients on $B_r$, and
extend the lower-order coefficients boundedly.  Equations
\eqref{eq:full-cauchy-data}--\eqref{eq:second-extension-derivative} show that
no distribution supported on $I_r$ occurs, so $\widetilde w$ satisfies an
inequality of the form \eqref{eq:ucp-inequality} in the full ball.  Since
$\widetilde w$ vanishes on the open lower half-ball,
Lemma~\ref{lem:ucp-external} yields $\widetilde w\equiv0$ in a smaller ball.
Thus the two prescribed sheets agree near the chosen boundary point.

It remains to patch the local identifications.  Cover the compact boundary
component by finitely many sheetwise embedding neighborhoods.  On an overlap,
the two local identifications both fix the boundary parametrization and map
one chosen embedded sheet onto the same chosen embedded sheet.  After
shrinking the cover, local injectivity of either immersion forces the two
identifications to coincide on the overlap.  They therefore patch to a
one-sided collar diffeomorphism.  Finally compactness permits one common positive
collar width, which completes the proof.
\end{proof}

\end{document}
